\documentclass[12 pt, letterpaper]{hmcpset}

\usepackage{lin-alg-preamble}

\name{Jayden Thadani}
\class{MATH 171 --- Abstract Algebra I}
\assignment{Homework 2}
\duedate{27th January 2024}

\begin{document}

\begin{problem}[Dummit \& Foote \S 1.1 \#25]
	Let $G$ be a group. Prove that if $x^2 = 1$ for all $x \in G$, then $G$ is abelian.
\end{problem}

\begin{solution}
	For any two elements $g, h \in G$, and identity $1$, we must have $g^2 = h^2 = (gh)^2 = 1$. Thus, $g, h, gh$ are all self-inverses. \\\\
	Thus,
	\[
		\begin{split}
			g h & = (g h)^{-1} \\
			    & = h^{-1} g^{-1} \\
			    & = h g
		\end{split}
	\]
	That is, any two arbitrary elements commute and then, by definition, $G$ is abelian.
\end{solution}

\pagebreak

In exercises \#2 and \#3, $D_{2n}$ has the usual presentation $D_{2n} = \left < r, s \vert r^n = s^2 = 1, rs = sr^{-1} \right >$

\begin{problem}[Dummit \& Foote \S  1.2 \#2]
	Use the generators and relations above to show that if $x$ is any element of $D_{2n}$ which is not a power of $r$, then $rx = xr^{-1}$.
\end{problem}

\begin{solution}
	Note that we have every element of $D_{2n}$ as $r^k$ or $s r^k$ for some integer $k$. Thus, all of the elements that are not powers of $r$ are of the form $s r^k$ with $k \in \mathbb Z$. Thus, for $x \in D_{2n}$, $x$ not a power of $r$, we must have
	\[
		\begin{split}
			r x & = r s r^k \\
			    & = s r^{-1} r^k \\
			    & = s r^{k - 1} \\
			    & = s r^k r^{-1} \\
			    & = x r^{-1}.
		\end{split}
	\]
	as desired. Note that the intermediate second and third steps hold because of associativity.
\end{solution}

\pagebreak

\begin{problem}[Dummit \& Foote \S 1.2 \#3]
	Use the generators and relations above to show that every element of $D_{2n}$ which is not a power of $r$ has order $2$. Deduce that $D_{2n}$ is generated by the two elements $s$ and $sr$, both of which have order $2$.
\end{problem}

\begin{solution}
	Again, we use the fact that all of the elements of $D_{2n}$ that are not a power of $r$ are of the form $s r^k$ for some integer $k$.\\\\ 
	Since we have $s \neq r^m$ for any integer $m$, $s r^k \neq 1 = r^n$. That is, it doesn't have order $1$. In order to show that it has order two, we will first show by induction that we can generalise the relation $r s = s r^{-1}$ to get $r^m s = s r^{-m}$ for $x$. \\\\
	In the case where $m = 1$, we already have $r^m s = s r^{-m}$. Suppose we had $r^m s = s r^m$ for some $m \in \mathbb N$. Then we would have
	\[
		\begin{split}
			r^{m + 1} s & = r r^m s \\
				    & = r s r^{-m} \\
				    & = s r^{-1} r^{-m} \\
				    & = s r^{-(m + 1)}.
		\end{split}
	\]
	Thus, by induction, we have $r^m s = s r^{-m}$ for all $m \in \mathbb N$. \\\\
	This lemma then gives us
	\[
		\begin{split}
			(s r^k)^2 & = s r^k s r^k \\
				  & = s s (r^k)^{-1} r^k \\
				  & = 1.
		\end{split}
	\]
	Thus, any $x \in D_{2n}$ not a power of $r$ has order $2$.
\end{solution}

\pagebreak

\begin{problem}[Dummit \& Foote \S 1.2 \#4]
	If $n = 2k$ is even and $n \ge 4$, show that $z = r^k$ is an element of order $2$ which commutes with all elements of $D_{2n}$. Show also that $z$ is the only nonidentity element of $D_{2n}$ which commutes with all elements of $D_{2n}$.
\end{problem}

\begin{solution}
	We can show that $r^k$ has order $2$, by seeing that $(r^k)^2 = r^{2k} = r^n = 1$. Now we show that $r^k$ commutes with all of $D_{2n}$\\\\
	Suppose $g = s^i r^j$ be an arbitrary element in $D_{2n}$ (with $i = \{ 0, 1 \}$, $j \in \{ 0, \ldots, n - 1 \}$). Then, if $i = 0$ consider
	\[
		\begin{split}
			r^k g & = r^k r^j \\
			      & = r^{k + j} \\
			      & = r^j r^k \\
			      & = g r^k
		\end{split}
	\]
	and if $i = 1$
	\[
		\begin{split}
			r^k g & = r^k s r^j \\
			      & = s r^{-k} r^j \\
			      & = s r^{j - k} \\
			      & = s r^j r^{-k} \\
			      & = g r^k
		\end{split}
	\]
	where the last step holds because $r^k$ having order $2$ implies that $r^k = r^{-k}$. \\\\
	Now we will show that $r^k$ is the only nonidentity element that commutes with all elements. Suppose some fixed $h = s^l r^m \in D_{2n}$ commutes with all elements $g = s^i r^j \in D_{2n}$ (with $l, i = \{ 0, 1 \}$, $m, j \in \{ 0, \ldots, n - 1 \}$). Then we have
	\[
		s^l r^m s^i r^j = s^i r^j s^l r^m
	\]
	and thus,
	\[
		s^{l + i} r^{j - m} = s^{l + i} r^{j + m}.
	\]
	This finally gives us
	\[
		r^{j - m} = r^{j + m}.
	\]
	Since $\lvert r \rvert = n$, we can reduce this to
	\[
		j - m = j + m \pmod n
	\]
	which then finally gives us
	\[
		2m = n \text{ or } m = 0.
	\] \\
	Thus, we have that $h$ must be $s^i r^k$ for some $i \in \{ 0, 1 \}$ (we ignore $m = 0$ since that's just the identity or $s$ and $s$ trivially does not commute with $r$). We show that $sr^k$ doesn't commute with $r$ and then we're done ---
	\[
		\begin{split} 
			r s r^k & = s r^{k - 1} \\
				& \neq s r^{k + 1} = s r ^k r
		\end{split}
	\] 
	for $n \ge 4$ and $k \ge 2$.
	Thus, $h = r^k$ is the only nonidentity element of $D_{2n}$ that commutes with all of $D_{2n}$.
	
\end{solution}

\pagebreak

\begin{problem}[Dummit \& Foote \S 1.2 \#5]
	If $n$ is odd and $n \ge 3$, show that the identity is the only element of $D_{2n}$ which commutes with all the elements of $D_{2n}$.
\end{problem}

\pagebreak

\begin{problem}[Dummit \& Foote \S 1.1 \#22 (optional)]
	If $x$ and $g$ are elements of the group $G$, prove that $\lvert x \rvert = \vert g^{-1} x g \rvert$. Deduce that $\lvert a b \rvert = \lvert b a \rvert$ for all $a, b \in G$.
\end{problem}

\begin{solution}
	First we will show that $(g^{-1} x g)^n = g^{-1} x^n g$. We will do so by induction. \\\\
	Suppose $x^m = 1$, for $m = 1$. Then
	\[
		\begin{split}
			(g^{-1} x g)^m & = g^{-1} x g \\
				       & = g^{-1} x^m g
		\end{split}
	\] \\
	Assume, that for $m \in \mathbb N$, we have that $(g^{-1} x g)^m = g^{-1} x^m g$. Then,
	\[
		\begin{split}
			(g^{-1} x g)^{m + 1} & =  g^{-1} x^m g g^{-1} x g \\
					     & = g^{-1} x^{m + 1} g.
		\end{split}
	\]
	Thus, for all $n \in \mathbb N$, we have $(g^{-1} x g)^n = g^{-1} x^n g$. \\\\
	Now we can show that $x^n = 1$ if and only if $(g^{-1} x g)^n = 1$, which we will call \textit{Lemma A} \\\\
	Suppose $x^n = 1$. Then
	\[
		\begin{split}
			(g^{-1} x g)^n & = g^{-1} x^n g \\
				       & = g^{-1} 1 g \\
				       & = g^{-1} g \\
				       & = 1.
		\end{split}
	\] \\
	Suppose $(g^{-1} x g)^n = 1$. Let $y = (g^{-1} x g)^n$ and $h = g^{-1}$. Then, since $y^n = 1$, we have $(h^{-1} y h)^n = 1$. But
	\[
		\begin{split}
			(h^{-1} y h)^n & = h^{-1} y^n h \\
				       & = g (g^{-1} x g)^n g^{-1} \\
				       & = g g^{-1} x^n g g^{-1} \\
				       & = x^n
		\end{split}
	\]
	and thus, $x^n = 1$. \\\\
	Now, finally, we can show that $\lvert x \rvert = \lvert g^{-1} x g \rvert$. Suppose by way of contradiction that $\lvert x \rvert \neq \lvert g^{-1} x g \rvert$. Suppose, without loss of generality, $\lvert x \rvert < \lvert g^{-1} x g \rvert$. If it's the former, then by \textit{Lemma A}
	\[
		(g^{-1} x g)^{\lvert x \rvert} = 1
	\]
	which contradicts that $\lvert g^{-1} x g \rvert$ is the least power of $g^{-1} x g$ that gives $1$. Thus, our hypothesis must be false, and $\lvert x \rvert = \lvert g^{-1} x g \rvert$
\end{solution}

\pagebreak

\begin{problem}[Dummit \& Foote \S 1.1 \#27 (optional)]
	Prove that if $x$ is an element of the group $G$ then $\{ x^n | n \in \mathbb Z \}$ is a subgroup of $G$ (called the \textit{cyclic subgroup} of $G$ generated by $x$).
\end{problem}

\begin{solution}
	Let $H = \{ x^n | n \in \mathbb Z \}$. For any $h_1 = x^m$ and $h_2 = x^n$ in $H$, we have
	\[
		\begin{split}
			h_1 h_2 & = x^n x^m \\
				& = x^{n + m} \in H
		\end{split}
	\]
	and
	\[
		\begin{split}
			h_1^{-1} & = (x^n)^{-1} \\
				 & = x^{-n} \in H.
		\end{split}
	\]
	Thus, $H$ is a subgroup, by the definition given in problem \#26.
\end{solution}

\end{document}

